% Hafta 14 — Tigress dönüşümleri ↔ 9. hafta kuralları
% Derleme: py -3.12 tools/latex/derle.py docs/week-14/assets/h14-03-donusum-esleme.tex
\documentclass[tikz,border=8pt]{standalone}
\usepackage{cen429-cizim}
\begin{document}
\begin{tikzpicture}[node distance=5mm and 8mm]
  \node[baslik] (b) at (0,0) {Tigress dönüşümleri $\leftrightarrow$ 9. hafta kuralları};
  \node[altbaslik, text width=150mm] at (b.south west) {araç yeni bir şey icat etmiyor --- aynı kuralları otomatikleştiriyor};

  \node[kutu=58mm, below=14mm of b.south west, anchor=north west] (k1) {\kb{Flatten}};
  \node[etiket, right=6mm of k1, anchor=west, text width=80mm] {K-04 kontrol akışı düzleştirme};
  \node[kutu=58mm, below=of k1] (k2) {\kb{AddOpaque / InitOpaque}};
  \node[etiket, right=6mm of k2, anchor=west, text width=80mm] {K-01 opak yüklem · K-03 ölü dal};
  \node[kutu=58mm, below=of k2] (k3) {\kb{EncodeArithmetic}};
  \node[etiket, right=6mm of k3, anchor=west, text width=80mm] {K-02 aritmetik kodlama};
  \node[kutu=58mm, below=of k3] (k4) {\kb{EncodeLiterals}};
  \node[etiket, right=6mm of k4, anchor=west, text width=80mm] {K-07 sabit / dize gizleme};
  \node[koyukutu=58mm, below=of k4] (k5) {\kb{Virtualize}};
  \node[etiket, right=6mm of k5, anchor=west, text width=80mm] {K-10 sanallaştırma (en pahalı)};

  \node[dip, text width=160mm, anchor=north west] at ($(k1.west |- k5.south)+(0,-8mm)$) {%
    Hangi dönüşümü seçtiğinizi S9'da kuralla eşleyerek gerekçelendirin.};
\end{tikzpicture}
\end{document}
